% Entwurf Abschnitt 6.1 (AP3, Plan review/plan-kap6.md Abschnitt 5). Stand 20.11.2026.
% Umfang Ziel ca. 1,7 S. (6.1.1 ca. 0,3 / 6.1.2 ca. 0,6 / 6.1.3 ca. 0,8).
% Code-Stand: G1@25d5c1a; übernommener Stand = G1@6c0f331 ohne die Simulationsdateien
% (docker-compose.sim.yaml, docker/Dockerfile.sim, docker/Dockerfile.mujoco, launch/bringup_sim.launch.py,
% G1_SIM_MODE-Zweige), die laut Nils (N3, 16.11.2026) Eigenleistung sind.
% Der Upstream-Stand 455e287/efc5e7c (hucebot/g1pilot) liegt nicht im Repo; die Defizite sind am Code von 6c0f331 geprüft.

\section{Ausgangsbasis und Abgrenzung}
\label{sec:ausgangsbasis}
% Basis: G1@25d5c1a g1pilot/NOTICE.md, Abschnitt „g1pilot (Basis dieses Pakets)“
Die Umgebung baut auf dem quelloffenen ROS-2-Paket G1Pilot auf.
Dieser Abschnitt begründet diese Wahl, gleicht den übernommenen Funktionsumfang mit den Anforderungen A1 bis A7 ab und grenzt die Eigenleistung ab.

\subsection{Auswahl der Ausgangsbasis}
\label{sec:auswahl_basis}
% Basis: Nils, 16.11.2026 (N4 in review/plan-kap6.md): einzige echte Alternative war die Neuentwicklung, g1pilot passte zum Vorhaben.
% Armregler mit IK: G1@6c0f331 g1pilot/g1pilot/manipulation/arm_controller.py (Klasse ArmController), g1pilot/g1pilot/utils/ik_solver.py.
% Lizenz: G1@25d5c1a g1pilot/LICENSE und g1pilot/NOTICE.md (BSD-3-Clause).
Als Ausgangsbasis stehen eine Neuentwicklung oder einer der Softwarestände aus Abschnitt~\ref{sec:verwandte_arbeiten} zur Wahl, die Tabelle~\ref{tab:stacks} gegenüberstellt.
Die Wahl fällt auf G1Pilot, weil dessen vorhandene Bausteine zum Vorhaben passen und eine Neuentwicklung sie erneut erfordert.
G1Pilot ist ein quelloffenes ROS-2-Paket für den G1, das den Unterkörper dem Unitree-Regler überlässt und Werkzeuge für die Steuerung des Oberkörpers und die Teleoperation bereitstellt \cite{HucebotG1Pilot}\ev{HucebotG1Pilot.01}.
Der übernommene Stand enthält zudem einen Armregler für kartesische Zielvorgaben auf Basis inverser Kinematik und steht unter der BSD-3-Clause-Lizenz.
Den Preis der Entscheidung benennt Abschnitt~\ref{sec:defizite}.

\subsection{Übernommener Funktionsumfang und Defizite}
\label{sec:defizite}
% Basis: eigene Analyse G1@6c0f331 (g1pilot/README.md „Nodes Overview“; g1pilot/g1pilot/manipulation/, navigation/,
% state/, teleoperation/, utils/) und G1@25d5c1a g1pilot/NOTICE.md (Tabelle Original-Anteil, Satz „Entfernt wurden …“).
Aus G1Pilot übernimmt die Umgebung sechs Bausteine: das Grundgerüst der Armregelung mit dem Löser der inversen Kinematik, die Zielmarker für Handposen in RViz, die Anbindung des Onboard-Reglers für Stehen und Gehen, die Veröffentlichung des Roboterzustands, die Steuerung über einen Joystick und einen Pfadplaner für die Navigation.
Nicht übernommen sind die Ansteuerung der Dex3-Hand, die tastenbasierte Bedienoberfläche, die Anbindung von LiDAR und Odometrie sowie das Docker-Image des Originals.

% Basis: eigene Analyse G1@6c0f331:
%  A2: navigation/loco_client.py (Stehen/Gehen nur über LocoClient; use_robot:=false -> self.robot = None, entering_balancing meldet nur „Sim balancing done“);
%      manipulation/arm_controller.py Docstring Klasse + main_loop-Ende (ohne Roboter nur JointState-Ausgabe); N3 (Nils): Original nur für Realbetrieb.
%  A5: arm_controller.py main_loop (kNotUsedJoint0.q = 1.0 ab dem ersten Zyklus, keine Rampe; bei arms_enabled=False kein Write mehr);
%      loco_client.py emergency_callback und Taste 5 (Damp() des Onboard-Reglers, Arme frei); keine Kollisionsprüfung (grep „collision“ ohne Treffer);
%      Begrenzung nur über Parameter arm_velocity_limit bzw. config.yaml inverse_kinematics.velocity_limit.
%  A6: teleoperation/ui_interface.py (Schaltflächenraster 5x5, nur Publisher, keine Subscriber -> keine Zustandsanzeige);
%      README.md „Instructions“ (Start über colcon build, ros2 launch).
%  A1/A7: keine HTTP-/Socket-Schnittstelle (grep ohne Treffer); Posen nur Grundstellung „home“ in config/config.yaml;
%      Hände nur manipulation/dx3_hand.py.
Der Abgleich mit den Anforderungen aus Tabelle~\ref{tab:anforderungen} zeigt vier Defizite des übernommenen Stands, die eine eigene Analyse des Codes belegt:
\begin{itemize}
  \item \textbf{A2:} Am schwersten wiegt, dass G1Pilot auf den Realbetrieb ausgelegt ist und keine physikalische Simulation enthält; jede Funktion lässt sich so nur am realen Roboter erproben. Ohne Roboter gibt der Armregler die Sollwinkel nur zur Anzeige aus, Stehen und Gehen sind allein über den Onboard-Regler erreichbar.
  \item \textbf{A5:} Der Armregler übernimmt die Arme ohne Übergang des Gewichtswerts der Armschnittstelle (Abschnitt~\ref{sec:g1_armsdk}) und gibt sie ohne Rückführung ab. Der Not-Halt ist die einzige Haltstufe und dämpft über den Onboard-Regler auch die Beine. Eine Kollisionsprüfung fehlt; die Armbewegung begrenzt allein ein einstellbarer Höchstwert der Gelenkgeschwindigkeit.
  \item \textbf{A6:} Die Bedienoberfläche bietet alle Funktionen gleichzeitig als Schaltflächenraster an und zeigt keinen Zustand des Roboters an. Der Start setzt Kenntnisse der ROS-2-Kommandozeile voraus.
  \item \textbf{A1 und A7:} Für Programme ohne ROS~2 fehlt eine Schnittstelle. Außer einer fest eingestellten Grundstellung lassen sich keine Posen oder Abläufe speichern, und angebunden ist nur die Dex3-Hand statt der Inspire-Hände des eingesetzten Roboters (Abschnitt~\ref{sec:g1_hardware}).
\end{itemize}
Aus den vier Defiziten folgen die Handlungsfelder der Abschnitte~\ref{sec:architektur} bis~\ref{sec:erweiterungskonzept}: Architektur mit Simulationspfad, Sicherheit, Bedienung und Erweiterung.

\subsection{Abgrenzung der Eigenleistung}
\label{sec:eigenleistung}
% Basis: G1@25d5c1a g1pilot/NOTICE.md (Original-Anteil je Datei; „etwa ein Zehntel der Zeilen“, dort Stand 30.10.2026);
% Stufen-Zuordnung nach Plan Abschnitt 5, 6.1.3 und tab:eigenleistung (AP2).
% Eigene Zählung 20.11.2026 an G1@25d5c1a (nicht im Text, bei Bedarf ergänzen): 26.398 Zeilen in g1pilot/g1pilot/**/*.py,
% 13.367 Zeilen in g1pilot/test/**/*.py, 597 Funktionen „def test_“ in g1pilot/test (deckt sich mit der Zählung vom 16.11. im Plan).
Die Abgrenzung der Eigenleistung folgt Funktionsbereichen, nicht Codezeilen.
Nach eigener Auswertung stammt zwar nur noch etwa ein Zehntel der Zeilen aus dem Original, viele Dateien mit Original-Code enthalten aber überwiegend eigene Erweiterungen, sodass eine Zeilenbilanz die Herkunft einzelner Funktionen verdeckt.
Tabelle~\ref{tab:eigenleistung} ordnet deshalb jedem Funktionsbereich eine von drei Stufen zu: \emph{übernommen} (aus G1Pilot im Kern unverändert), \emph{erweitert} (Grundgerüst aus G1Pilot, wesentliche Teile neu) oder \emph{neu} (in G1Pilot nicht enthalten).

% Basis: G1@25d5c1a g1pilot/NOTICE.md (robot_state.py ~80 %; dijkstra_planner.py, nav2point.py, create_map.py 35–80 %;
% arm_controller.py Grundgerüst ~20 %; ik_solver.py ~30 %; loco_client.py Grundgerüst ~20 %;
% interactive_marker.py ~50 %, joystick.py/joy_mux.py 25–80 %); docs/13-A-Alter-Joystick.md „Richtet sich an: Entwickler und Betreuer“.
Übernommen sind die Veröffentlichung des Roboterzustands und der Pfadplaner, den die Umgebung für die Navigation in der Simulation einbindet; die Zielmarker in RViz und der Joystick-Zugang bleiben als Werkzeuge für Entwickler erhalten.
Erweitert sind die Armregelung mit der inversen Kinematik und die Anbindung des Onboard-Reglers, deren Grundgerüst aus G1Pilot um die Schutz- und Übergabefunktionen aus Abschnitt~\ref{sec:sicherheitskonzept} ergänzt ist.

% Basis: G1@25d5c1a g1pilot/NOTICE.md, Satz „Alle übrigen Dateien (u. a. Demo-GUI, Arm-Planer, loco_sim mit Lauf-Policies,
% Inspire-Hände, PS4/Xbox-Teleop, Arm-API, Handy-Zugang, Tests, Doku) sind eigene Arbeit“; N3 (Simulation vollständig Eigenleistung).
Den größten Teil der Tabelle bilden neue Bereiche: das Kollisionsmodell mit dem Bahnplaner, die Kopplung an die Simulation, der Ersatzregler für Stehen und Gehen und die Sicherheitsfunktionen.
Neu sind ebenso die Inspire-Hände mit eigenen Hand-Profilen, die Bedienoberfläche, die Gamepad-Bedienung und der Smartphone-Zugang, die Posen und Abläufe, die Schnittstelle für externe Programme, die Bereitstellung mit Startmenü sowie Tests und Dokumentation.

% Basis: G1@25d5c1a g1pilot/NOTICE.md, Tabelle „Weitere Fremdbestandteile“; g1pilot/policies/g1_wholebody/ATTRIBUTION.md,
% g1pilot/policies/agile_velocity_g1/ATTRIBUTION.md, g1pilot/policies/agile_velocity_height_g1/ATTRIBUTION.md;
% unitree_mujoco/scene_editor/README.md (Grundlage mujoco-scene-editor).
Nicht zur Eigenleistung zählen die eingebundenen Fremdbestandteile.
Die Kopplung an die Simulation baut auf unitree\_mujoco auf, einem Simulator des Herstellers auf Basis von unitree\_sdk2 und MuJoCo \cite{UnitreeMujoco}\ev{UnitreeMujoco.01}.
Hinzu kommen weitere Herstellerpakete, ein browserbasierter Szenen-Editor für MuJoCo und die Lauf-Policies aus Unitree RL Mjlab, einem Lernprojekt mit MuJoCo als Physik-Backend \cite{UnitreeRLMjlab}\ev{UnitreeRLMjlab.01}, und aus NVIDIA WBC-AGILE \qf{WBC-AGILE: Zhao2026, Freigabe fehlt}.
Abschnitt~\ref{sec:eval_code} greift die Abgrenzung auf, wenn er den Programmstand von Simulation und Realbetrieb vergleicht.
